\documentclass{article}
\usepackage{graphicx} % Required for inserting images

\title{STT-2902}
\author{Julien Miron}


\begin{document}

\section*{Équipe 2 – Différences entre les types d’échantillonnage}

Objectif : Expliquer différents types d’échantillonnage (aléatoire simple, systématique, stratifié, par grappes). Votre présentation ne traite pas des biais directement, mais comme on choisit souvent le type d'échantillonnage pour réduire ceux-ci, il est nécessaire d'en parler.

\vspace{1cm}
Suggestions :
\begin{itemize}
    \item Proposez une activité interactive : chaque équipe doit simuler un échantillonnage différent à partir d’une "population" fictive (boules de couleurs, cartes, etc.).
    \item Discutez ensuite des biais potentiels selon le type choisi.

\end{itemize}

\vspace{1cm}
\textbf{Requis} :
\begin{enumerate}
    \item Illustrer, avec un exemple concret, \textbf{un biais} causé par un mauvais choix d’échantillonnage.
    \item Fournir un schéma ou diagramme qui résume visuellement les différences entre les types d’échantillonnage.
\end{enumerate}


\end{document}
